\subsection{\texttt{TROWARGMIN}}
\noindent\textbf{Bundle encoding.} UOP entry 43, opcode \texttt{0xF001}. Bundle: \texttt{B.OP + B.ARG + B.DIM + B.IOTI}. \texttt{B.OP.Opcode}=\texttt{0xF001}. \texttt{B.IOTI}: selector=\texttt{111 last}, \texttt{SrcTile0}=\texttt{src}, \texttt{SrcTile1}=\texttt{absent}, \texttt{DstTile}=\texttt{dst}, \texttt{SizeCode}=profile tile size. \texttt{B.DIM} supplies row-axis reduction or row-broadcast bounds. Scratch operand(s) \texttt{tmp} are implementation-local and are not encoded in the CPU-visible bundle.\par\smallskip
{\small
\paragraph{PTO ISA description.}
Get the column index of the minimum element for each row.
\paragraph{Example.}
Source row [3, 1, 7, 2]: output row[0] = 1 (0-indexed column of min value 1).
\par\smallskip
\paragraph{PTO ISA operation.}
Let \texttt{R = src.GetValidRow()} and \texttt{C = src.GetValidCol()}. For \texttt{0 \textless{}= i \textless{} R}:

\[
\mathrm{dst}_{i,0} = \operatorname*{arg\,min}_{0 \le j < C}\; \mathrm{src}_{i,j}
\]

\noindent where the result is the column index $j$ that minimises $\mathrm{src}_{i,j}$.
}\par\smallskip
\paragraph{Notes.}
Same tie-breaking as \texttt{TROWARGMAX}. Output type should be an integer type.
\paragraph{Microcode site table.}
{\scriptsize
\begin{longtable}{@{}R{0.055\linewidth}L{0.23\linewidth}L{0.31\linewidth}L{0.22\linewidth}@{}}
\toprule
Seq & Micro instruction & WO[63:32] fields & WM[31:0] fields \\
\midrule
0 & \texttt{u\_vec.vbr} & \texttt{opc=0x0B, x=0, fl=alu, seq=0, kind=c} & \texttt{entry=43, cnt=0, res=vec, var=0} \\
1 & \texttt{u\_vec.vbr} & \texttt{opc=0x0B, x=0, fl=alu, seq=1, kind=c} & \texttt{entry=43, cnt=0, res=vec, var=0} \\
2 & \texttt{u\_vec.vlds} & \texttt{opc=0x17, x=0, fl=mem, seq=2, kind=b} & \texttt{entry=43, cnt=2, res=vec, var=0} \\
3 & \texttt{u\_vec.vcmin} & \texttt{opc=0x0E, x=0, fl=alu, seq=3, kind=b} & \texttt{entry=43, cnt=2, res=vec, var=0} \\
4 & \texttt{u\_vec.vdintlv} & \texttt{opc=0x11, x=0, fl=alu, seq=4, kind=u} & \texttt{entry=43, cnt=1, res=vec, var=0} \\
5 & \texttt{u\_vec.vbr} & \texttt{opc=0x0B, x=0, fl=alu, seq=5, kind=c} & \texttt{entry=43, cnt=0, res=vec, var=0} \\
6 & \texttt{u\_vec.vbitcast} & \texttt{opc=0x0A, x=0, fl=alu, seq=6, kind=u} & \texttt{entry=43, cnt=1, res=vec, var=0} \\
7 & \texttt{u\_vec.vadds} & \texttt{opc=0x08, x=0, fl=alu, seq=7, kind=b} & \texttt{entry=43, cnt=2, res=vec, var=0} \\
8 & \texttt{u\_vec.vcmp} & \texttt{opc=0x0F, x=0, fl=alu, seq=8, kind=b} & \texttt{entry=43, cnt=2, res=vec, var=0} \\
9 & \texttt{u\_vec.vsel} & \texttt{opc=0x24, x=0, fl=alu, seq=9, kind=t} & \texttt{entry=43, cnt=3, res=vec, var=0} \\
10 & \texttt{u\_vec.pbitcast} & \texttt{opc=0x01, x=0, fl=setup, seq=10, kind=u} & \texttt{entry=43, cnt=1, res=vec, var=0} \\
11 & \texttt{u\_vec.vsel} & \texttt{opc=0x24, x=0, fl=alu, seq=11, kind=t} & \texttt{entry=43, cnt=3, res=vec, var=0} \\
12 & \texttt{u\_vec.vsts} & \texttt{opc=0x2A, x=0, fl=mem, seq=12, kind=b} & \texttt{entry=43, cnt=2, res=vec, var=0} \\
\bottomrule
\end{longtable}
}
\paragraph{Microcode body DSL.}
\begin{lstlisting}[style=ptocode]
def uop_trowargmin(src, tmp, dst):
    src_dtype = src.element_type
    idx_dtype = dst.element_type
    lanes = pto.get_lanes(src_dtype)
    valid_rows, valid_cols = src.valid_shape
    elem_bytes = pto.bytewidth(idx_dtype)

    init_val = pto.PadValue.MAX.eval(src_dtype)

    if pto.constexpr(elem_bytes == 4):
        store_dist = pto.VStoreDist.ONE_POINT_B32
    elif pto.constexpr(elem_bytes == 2):
        store_dist = pto.VStoreDist.ONE_POINT_B16
    else:
        store_dist = pto.VStoreDist.ONE_POINT_B8

    for row in range(0, valid_rows, 1):
        remained = valid_cols

        v_val_acc = pto.vbr(init_val)
        init_zero_idx = idx_dtype(0)
        v_idx_acc = pto.vbr(init_zero_idx)

        mask_1, _ = pto.make_mask(src_dtype, 1)
        mask_1_idx, _ = pto.make_mask(idx_dtype, 1)

        for col in range(0, valid_cols, lanes):
            mask, remained = pto.make_mask(src_dtype, remained)
            v_src = pto.vlds(src[row, col:])
            v_reduced = pto.vcmin(v_src, mask)

            v_val, v_idx = pto.vdintlv(v_reduced, pto.vbr(src_dtype(0)))
            v_idx = pto.vbitcast(v_idx, idx_dtype)

            col_offset = idx_dtype(col)
            v_idx = pto.vadds(v_idx, col_offset, mask_1_idx)

            cmp_mask = pto.vcmp(v_val_acc, v_val, mask_1, "gt")

            v_val_acc = pto.vsel(v_val, v_val_acc, cmp_mask)
            cmp_mask_b32 = pto.pbitcast(cmp_mask, pto.mask_b32)
            v_idx_acc = pto.vsel(v_idx, v_idx_acc, cmp_mask_b32)

        pto.vsts(v_idx_acc, dst[row, 0:], mask_1_idx, dist=store_dist)
    return
\end{lstlisting}
\Needspace{0.34\textheight}
